% =====================================================================
% Routing cluster maps: original-order belong_idx reshaped to HxW.
% Two-column figure (figure*); \input into the paper after confidence histogram.
%
% ASSET PROVENANCE (paper/figures/fig58_assets/, generated on the training
%   machine by CATANet/scripts/build_routing_figures.py with the x4 model
%   net_g_250000): for each hard sample we forward the LR image once and color
%   each LR pixel by its prototype index b_n=argmax_m S_{n,m} at three TAB
%   blocks (b0: M=16, b3: M=128, b7: M=128). Colors are a fixed golden-ratio
%   hue palette (distinct hue per prototype); the same palette indexes all
%   panels but prototype identities are per-block (not comparable across blocks).
%   "LR" is the nearest-neighbor-magnified input for reference. Hard samples are
%   the same as Fig.7 (Urban100 img_092/img_024 dense grids; Manga109 line art).
% =====================================================================
\newcommand{\figeightheader}{%
  \begin{minipage}{0.245\textwidth}\centering{\scriptsize Reference LR}\end{minipage}\hfill
  \begin{minipage}{0.245\textwidth}\centering{\scriptsize Early router: coarse groups}\end{minipage}\hfill
  \begin{minipage}{0.245\textwidth}\centering{\scriptsize Middle router: finer groups}\end{minipage}\hfill
  \begin{minipage}{0.245\textwidth}\centering{\scriptsize Last router: fine groups}\end{minipage}%
}

\newcommand{\figeightrow}[2]{%
  % #1 = asset tag   #2 = row caption (shown under the LR panel)
    \setlength{\fboxsep}{0pt}%
    \setlength{\fboxrule}{0.25pt}%
  \begin{minipage}{0.245\textwidth}\centering
      \fbox{\includegraphics[width=0.985\linewidth]{figures/fig58_assets/fig8_#1_lr.png}}\\[1pt]
    {\tiny #2}
  \end{minipage}\hfill
  \begin{minipage}{0.245\textwidth}\centering
      \fbox{\includegraphics[width=0.985\linewidth]{figures/fig58_assets/fig8_#1_b0_cluster.png}}\\[1pt]
      {\tiny $b_0$ coarse ($M{=}16$)}
  \end{minipage}\hfill
  \begin{minipage}{0.245\textwidth}\centering
      \fbox{\includegraphics[width=0.985\linewidth]{figures/fig58_assets/fig8_#1_b3_cluster.png}}\\[1pt]
      {\tiny $b_3$ finer ($M{=}128$)}
  \end{minipage}\hfill
  \begin{minipage}{0.245\textwidth}\centering
      \fbox{\includegraphics[width=0.985\linewidth]{figures/fig58_assets/fig8_#1_b7_cluster.png}}\\[1pt]
      {\tiny $b_7$ fine ($M{=}128$)}
  \end{minipage}%
}

\begin{figure}[H]
\centering
\figeightheader
\vspace{2pt}

\figeightrow{Urban100_img_092}{Urban100: img\_092 (LR)}

\vspace{2pt}
\figeightrow{Urban100_img_024}{Urban100: img\_024 (LR)}

\vspace{2pt}
\figeightrow{Manga109_ThatsIzumiko_000}{Manga109: ThatsIzumiko (LR)}

\caption{Routing cluster maps at $\times4$. Each LR pixel is colored by its
assigned prototype index $b_n=\arg\max_m S_{n,m}$ (Eq.~\eqref{eq:belong}) and
reshaped from the original token order to $H\times W$. Read each row by comparing
the cluster boundaries with the reference LR image: the early block forms coarse
content groups, while the deeper blocks split repeated facade elements and
high-contrast strokes into finer partitions. Colors encode prototype identity
within a block and are not comparable across blocks, so the meaningful cue is
region continuity rather than a specific hue. This visualizes the same routing
signal whose confidence $x^{\mathrm{s}}$ is quantified in
Figure~\ref{fig:xscore_hist}.}
\label{fig:cluster_maps}
\end{figure}
